\documentclass{ieeeaccess}
\usepackage{cite}
\usepackage{url}
\usepackage{amsmath,amssymb,amsfonts}
\usepackage{hyperref}
\usepackage{tabularx}
\usepackage{graphicx}
\usepackage[numbers, sort]{natbib}
\usepackage{algorithm}
\usepackage{algpseudocode}
\usepackage{bm}
\usepackage{textcomp}
\usepackage{array}
\usepackage{caption}
\usepackage{booktabs}
\usepackage{times}
\usepackage{adjustbox}
\usepackage{multirow}
\usepackage{xcolor}
\usepackage{float}
\usepackage{booktabs,siunitx}
\usepackage{subcaption}

\newcolumntype{M}{>{\small\bfseries}l}

\def\BibTeX{{\rm B\kern-.05em{\sc i\kern-.025em b}\kern-.08em
T\kern-.1667em\lower.7ex\hbox{E}\kern-.125emX}}

\begin{document}

\history{Date of publication xxxx 00, 0000, date of current version xxxx 00, 0000.}
\doi{10.1109/ACCESS.2017.DOI}

\title{DataFlow~AI: A Hybrid, Confidence-Aware Framework for Enterprise Data-Integrity Validation and Cross-Dialect SQL Translation}

\author{
  \uppercase{Bhaumik Upasana}\authorrefmark{1},
  \uppercase{M~A~Berlin}\authorrefmark{2}
}

\address[1,2]{School of Computer Science and Engineering (SCOPE), Vellore Institute of Technology (VIT), Vellore, Tamil Nadu 632 014, India}

\corresp{Corresponding author: Berlin M A (e-mail: berlin.ma@vit.ac.in).}

\begin{abstract}
Enterprise data pipelines lose substantial value to poor data quality, and database-migration projects frequently overrun
or fail outright, in large part because integrity validation and query translation are treated as disconnected problems.
We present DataFlow~AI, a deployment-oriented framework that couples a hybrid statistical/machine-learning anomaly-detection
ensemble with a confidence-aware, cross-dialect SQL translation engine. The detector combines an Isolation Forest, TF--IDF
cosine similarity for near-duplicate identification, interquartile-range tests, and rule-based cross-field checks, merged
under a recall-maximising union strategy whose precision is governed post~hoc by a composite confidence router. The
translation engine uses Google Gemini~2.0~Flash with a deterministic rule-based fallback, attaching a calibrated confidence
score to every conversion. We evaluate the system on five synthetic financial datasets (totalling \num{20720} records with
known anomaly masks), a seven-point scalability sweep up to \num{100000} rows, five noise levels, twelve adversarial
security payloads, and a ten-statement SQL translation suite, all under a fixed random seed. Across four labelled datasets
the full ensemble attains a mean F1 of \num{0.32} at \SI{76}{\percent} recall, a \num{+0.19} absolute F1 gain over an
Isolation-Forest baseline that is statistically significant under \mbox{10-repeat $\times$ 5-fold} cross-validation
($t=51.2$, $p<0.0001$, Cohen's $d=9.1$). An ablation study isolates TF--IDF semantic matching as the dominant component
($\Delta\mathrm{F1}=-0.107$ on removal). The system scales sub-linearly (empirical exponent $\alpha=0.69$), handles all
adversarial and security inputs without failure, and translates SQL at \num{0.92} mean confidence with \num{0.97} mean
semantic similarity. We deliberately report the high false-positive regime that the union strategy produces and show how
confidence routing converts it into a practical, audit-appropriate review workflow. The results support hybrid,
confidence-aware integration over isolated single-method tools while remaining transparent about the operating point and
the limits of synthetic evaluation.
\end{abstract}

\begin{keywords}
Data integrity, anomaly detection, query translation, machine-learning ensemble, large language models, database
migration, confidence calibration, cross-field validation
\end{keywords}

\titlepgskip=-15pt

\maketitle

% =====================================================================
% TABLE 1 — abbreviations
\begin{table*}[tp]
\centering
\caption{List of Abbreviations}
\label{tab:abbreviations}
\begin{tabular}{|l|l||l|l|}
\hline
\textbf{Abbreviation} & \textbf{Definition} & \textbf{Abbreviation} & \textbf{Definition} \\
\hline
AI & Artificial Intelligence & API & Application Programming Interface \\
\hline
LLM & Large Language Model & DB & Database \\
\hline
TF-IDF & Term Frequency--Inverse Document Frequency & SQL & Structured Query Language \\
\hline
IF & Isolation Forest & ML & Machine Learning \\
\hline
RQ & Research Question & FPR & False Positive Rate \\
\hline
F1 & F1 Score (Harmonic Mean) & TP & True Positives \\
\hline
FP & False Positives & FN & False Negatives \\
\hline
TN & True Negatives & CI & Confidence Interval \\
\hline
CV & Cross-Validation & DDL & Data Definition Language \\
\hline
IQR & Interquartile Range & $\kappa$ & Cohen's Kappa \\
\hline
ROI & Return on Investment & QA & Quality Assurance \\
\hline
CPU & Central Processing Unit & RAM & Random-Access Memory \\
\hline
PostgreSQL & PostgreSQL Database & FastAPI & FastAPI Framework \\
\hline
React & React JavaScript Library & ORM & Object-Relational Mapping \\
\hline
DDL & Data Definition Language & DML & Data Manipulation Language \\
\hline
$O(n \log n)$ & Sub-linear/Log-linear Time Complexity & $\alpha$ & Scaling Exponent \\
\hline
XSS & Cross-Site Scripting & RTL & Right-To-Left (Unicode) \\
\hline
\end{tabular}
\end{table*}

% =====================================================================
\section{Introduction}
\label{sec:introduction}

\subsection{Data Quality and Migration as a Coupled Problem}
Modern organisations operate increasingly heterogeneous data estates, and two recurring failure modes dominate their
operational cost. The first is degraded data quality: industry surveys consistently report that data practitioners spend
the majority of their effort on preparation and cleaning rather than analysis, with poor quality cascading into downstream
analytics and compliance risk~\cite{redman1998impact,hultgren2018data}. The second is the difficulty of database
migration and modernisation: a large fraction of such projects overrun their schedules or are abandoned, and query
translation alone accounts for a substantial share of migration effort that remains stubbornly
manual~\cite{stonebraker2007end}.

These two problems are typically addressed by separate tools, yet in practice they are coupled. Migrating data of unknown
quality propagates errors into the target system, while validating data without a path to transform and load it into the
target schema leaves the value unrealised. Rule-based validators are difficult to configure per domain and produce
false-positive rates that erode trust in automation~\cite{batini2009data}, and most validators cannot detect semantic
inconsistencies that span columns, such as a name field that is inconsistent with the associated email or account
code~\cite{ilyas2015functional,motro1989integrity}. There remains a gap for integrated, deployment-aware systems that
treat validation and translation as parts of one workflow~\cite{stonebraker2007end}.

\subsection{Research Motivation}
Machine learning has transformed perception and language tasks, but data-engineering tooling remains fragmented: work
tends to optimise a single quality metric or a single schema-mapping step in isolation. We identify three gaps that a
practical framework must address.

\begin{enumerate}
\item \textbf{Operating-point opacity.} Detectors are frequently reported at a single, often favourable, threshold. A
deployable system must expose its precision/recall trade-off and provide a mechanism to manage it, because a high-recall
detector with an unmanaged false-positive rate is unusable in production.
\item \textbf{Absence of confidence estimation.} Translation tools rarely attach calibrated confidence to their output or
implement graceful fallbacks, forcing users to manually verify every converted query and negating the benefit of
automation~\cite{sculley2015hidden}.
\item \textbf{Sub-surface semantic violations.} Cross-column inconsistencies cannot be detected by univariate statistics
alone~\cite{motro1989integrity} and require explicit relational rules integrated into the same pipeline.
\end{enumerate}

\subsection{Research Questions}
\label{subsec:rqs}
We frame the evaluation around four research questions (RQs) that map directly to the gaps above and to the experiments
in Section~\ref{sec:results}.

\begin{enumerate}
\item \textbf{RQ1 (Detection quality).} Does the hybrid ensemble achieve materially higher anomaly-detection F1 and recall
than single-method baselines on diverse financial data, and is the improvement statistically significant under repeated
cross-validation?
\item \textbf{RQ2 (Component contribution).} Which components of the ensemble are responsible for its performance, and how
much does each contribute?
\item \textbf{RQ3 (Robustness, scalability, and security).} Does the system degrade gracefully under noise, scale
sub-linearly to large datasets, and handle adversarial and malformed inputs without failure?
\item \textbf{RQ4 (Translation and confidence).} Can LLM-assisted, fallback-protected translation convert SQL across
dialects with high semantic fidelity, and does the attached confidence behave as a usable routing signal?
\end{enumerate}

\subsection{Contributions}
DataFlow~AI is a hybrid, confidence-aware framework that integrates data-integrity validation and cross-dialect SQL
translation in a single deployment-oriented pipeline. Our contributions are:

\begin{enumerate}
\item A hybrid anomaly-detection ensemble (Isolation Forest, TF--IDF cosine similarity, IQR tests, and rule-based
cross-field checks) merged under a recall-maximising union strategy, whose precision is controlled by a composite
confidence router rather than by tuning each detector in isolation.
\item A confidence-aware translation engine that pairs an LLM with a deterministic rule-based fallback and attaches a
calibrated confidence score to every conversion, enabling a three-tier auto-accept / review / reject routing policy.
\item A rigorous, fully reproducible evaluation suite of nine experiments under a fixed seed, with an explicit and honest
account of the system's operating point: we report the high-FPR regime produced by the union strategy and demonstrate how
confidence routing renders it practical for audit contexts, rather than obscuring it behind a single favourable threshold.
\end{enumerate}

The remainder of the paper is organised as follows. Section~\ref{sec:related_work} reviews and critiques prior work.
Section~\ref{sec:architecture} describes the system architecture. Section~\ref{sec:methodology} details the detection and
translation methodology. Section~\ref{sec:setup} specifies the experimental setup and datasets.
Section~\ref{sec:results} reports results for RQ1--RQ4, including ablation, baseline, robustness, scalability, error, and
security analyses. Section~\ref{sec:casestudy} presents an enterprise case study. Section~\ref{sec:discussion} discusses
threats to validity, reproducibility, security considerations, limitations, and future work.
Section~\ref{sec:conclusion} concludes.

% =====================================================================
\section{Related Work}
\label{sec:related_work}

\subsection{Evolution of Data-Quality Validation}
Data-quality management has progressed through three broad phases. Early manual curation relied on expert review and did
not scale. Rule-based frameworks then automated common checks; tools in this tradition, exemplified by Great
Expectations~\cite{great_expectations}, codify expectations declaratively but require heavy per-domain configuration and
struggle with context-dependent semantics, emitting false alerts that must be confirmed by hand~\cite{batini2009data}.
More recent learning-based detectors~\cite{pedregosa2011scikit,aggarwal2017outlier} report encouraging accuracy on
isolated benchmarks but are seldom validated under production conditions, rarely integrate cross-column semantic checks,
and typically omit the confidence estimation and downstream integration that operational use demands.

A consistent and under-acknowledged theme across this literature is the dependence of reported metrics on the chosen
operating point. Detectors evaluated at a single threshold can appear excellent while concealing the precision/recall
trade-off that determines real-world reviewer burden. Our work treats that trade-off as a first-class object of study
rather than a tuning detail.

\subsection{Query and Schema Translation}
Commercial migration tooling has historically emphasised structural schema mapping over query translation, and industrial
reports indicate that query conversion still consumes a large share of migration effort even when partial automation is
available. Rule-based translators are deterministic and auditable but brittle on complex, nested, or vendor-specific
constructs~\cite{rahm2001survey}. Neural and few-shot approaches generalise better across dialects~\cite{manning2008introduction}
but introduce their own risks: they can hallucinate syntactically valid yet semantically incorrect output, and they rarely
quantify their own uncertainty or degrade gracefully when a service is unavailable. The community increasingly recognises
that production translation requires transparency, explainability, and calibrated confidence~\cite{ribeiro2016why}.

\subsection{Operational and Hybrid Machine Learning}
The difficulty of operating ML systems --- data drift, silent performance decay, and integration debt --- is well
documented~\cite{sculley2015hidden,neumann1996risks}. Hybrid designs that constrain probabilistic models with
deterministic rules can retain learning capacity while improving interpretability~\cite{lundberg2017unified}, and
operational frameworks additionally need confidence estimation and explicit fallback paths. DataFlow~AI adopts this hybrid
philosophy on both sides of its pipeline: rules and statistics bound the behaviour of the learned components in detection,
and a rule-based engine bounds the LLM in translation.

\subsection{Critical Synthesis and Positioning}
Synthesising across these threads, four limitations recur in prior systems: (i) single-threshold reporting that hides the
operating point; (ii) the inability to detect cross-field semantic violations with univariate methods; (iii) the lack of
calibrated confidence and graceful fallback in translation; and (iv) the treatment of validation and translation as
disjoint problems. DataFlow~AI is positioned precisely against these: it reports full confusion matrices and FPR rather
than a single headline number, integrates explicit cross-field rules, attaches confidence to every translation with a
deterministic fallback, and unifies detection and translation in one pipeline. We do not claim to set new accuracy
records; rather, we contribute a deployment-aware integration whose behaviour is fully characterised and reproducible.
Table~\ref{tab:comparative} summarises the functional coverage of representative systems against DataFlow~AI.

% TABLE 2 — comparative analysis (functional coverage; no inflated cross-system metric claims)
\begin{table*}[tp]
\centering
\caption{Functional Coverage of DataFlow~AI Relative to Representative Systems}
\label{tab:comparative}
\begin{tabular}{|l|c|c|c|c|c|c|}
\hline
\textbf{System} & \textbf{ML} & \textbf{Cross-Field} & \textbf{Confidence} & \textbf{Graceful} & \textbf{SQL} & \textbf{Unified} \\
 & \textbf{Anomaly} & \textbf{Semantic} & \textbf{Routing} & \textbf{Fallback} & \textbf{Translation} & \textbf{Pipeline} \\
\hline
Great Expectations~\cite{great_expectations} & No & No & No & No & No & No \\
\hline
Rule-based validators~\cite{batini2009data} & No & Partial & No & No & No & No \\
\hline
Learned detectors~\cite{aggarwal2017outlier} & Yes & No & No & No & No & No \\
\hline
Schema/query mappers~\cite{rahm2001survey} & No & No & No & Partial & Yes & No \\
\hline
Neural translators~\cite{manning2008introduction} & No & No & No & No & Yes & No \\
\hline
\textbf{DataFlow~AI (this work)} & \textbf{Yes} & \textbf{Yes} & \textbf{Yes} & \textbf{Yes} & \textbf{Yes} & \textbf{Yes} \\
\hline
\end{tabular}
\end{table*}

% =====================================================================
\section{System Architecture}
\label{sec:architecture}

\subsection{Overview and Technology Stack}
DataFlow~AI is implemented as a layered service with two functional modules --- a data-quality module and a migration
module --- sharing a common persistence and audit layer. The implementation uses FastAPI (Python~3.12+) for the
application layer, scikit-learn, pandas, NumPy, and SciPy for the analytical core, SQLAlchemy over SQLite for job state
and audit logging, and a React/Vite single-page application with WebSocket streaming for progress reporting. SQL
translation uses Google Gemini~2.0~Flash (temperature~0.1 for determinism) with a deterministic rule-based fallback.
Figure~\ref{fig:architecture} shows the three layers and the principal components within each.

\begin{figure*}[tp]
\centering
\includegraphics[width=0.82\textwidth]{figures/fig0_architecture.png}
\caption{DataFlow~AI layered architecture. The presentation layer streams live progress over WebSocket; the application
layer hosts the data-quality and migration modules, the hybrid detector, the confidence router, and the LLM translator
with rule fallback; the persistence layer provides ORM-backed state, an audit trail, and multi-format export. (New
architecture diagram, drawn from the implemented system.)}
\label{fig:architecture}
\end{figure*}

\subsection{Data-Quality Module}
The data-quality module exposes REST endpoints for assessing arbitrary tabular datasets. Its analytical core extracts
numeric features (standardised with a $z$-score scaler), concatenates text columns for vectorisation, and computes
per-row null rates. Six detectors run over these features: an Isolation Forest on standardised numerics, TF--IDF cosine
similarity for near-duplicate detection, IQR and $z$-score outlier tests, pattern validators (email, date, numeric
format), and rule-based cross-field consistency checks. A companion cleaner module provides imputation
($k$-nearest-neighbour and iterative), deduplication, outlier removal, and format standardisation; cleaning is out of
scope for the present evaluation, which focuses on detection and translation.

\subsection{Migration Module}
The migration module translates SQL between dialects (MySQL, PostgreSQL, and SQL~Server as sources; Snowflake and
PostgreSQL as targets). Its primary path is the LLM with structured prompting; its fallback path is a deterministic
rule-based converter covering structural and function-level transforms (for example, \texttt{IFNULL}$\rightarrow$\texttt{NVL},
backtick$\rightarrow$double-quote identifiers, \texttt{AUTO\_INCREMENT}$\rightarrow$\texttt{AUTOINCREMENT}, and removal of
engine/locking clauses). Enterprise features include batch scheduling with priority levels, multi-format export, and a
versioned migration history with rollback.

\subsection{Persistence and Real-Time Layer}
SQLite via SQLAlchemy stores migration jobs, per-statement translation records, quality jobs, per-issue findings, and an
append-only audit log. The WebSocket manager multiplexes per-job progress events (\texttt{progress},
\texttt{status\_update}, \texttt{error}, \texttt{complete}), enabling live monitoring of long-running detection and
translation jobs. The design targets single-instance deployment with a clear migration path to a server-grade RDBMS
through the ORM abstraction.

% =====================================================================
\section{Methodology}
\label{sec:methodology}

\subsection{Hybrid Detection Pipeline}
Figure~\ref{fig:pipeline} depicts the end-to-end detection pipeline. After ingestion and feature engineering, the three
principal detectors execute in parallel and their row-level flags are merged by a union (logical OR) rule, which maximises
recall. The merged flags are not the final decision: a composite confidence score (Section~\ref{subsec:confidence}) is
computed per row and used to route each record to one of three outcomes. This separation of a high-recall detector from a
precision-controlling router is the central design choice of the system, and it is the reason we report detector behaviour
and routing behaviour separately in Section~\ref{sec:results}.

\begin{figure}[tbp]
\centering
\includegraphics[width=\columnwidth]{figures/fig_pipeline.png}
\caption{Hybrid detection pipeline. Parallel detectors are merged under a recall-maximising union; a composite confidence
score then routes each row to auto-accept, manual review, or auto-reject. Routing percentages are the mean standard-threshold
rates measured across the financial datasets. (New diagram, drawn from the implemented pipeline.)}
\label{fig:pipeline}
\end{figure}

\subsubsection{Isolation Forest}
The Isolation Forest isolates anomalies through random recursive partitioning, exploiting the fact that outliers require
fewer splits to separate~\cite{liu2008isolation}. We use 100 trees, a contamination prior of $0.10$, and a fixed random
state of 42. The per-row anomaly score is normalised to $[0,1]$ and contributes both a binary flag and a continuous input
to the confidence score. On data with fewer than 20 rows the detector returns zeros, since no reliable fit is possible.

\subsubsection{TF--IDF Near-Duplicate Detection}
Text columns are concatenated and vectorised with TF--IDF (maximum 500 features, English stop-words removed), and per-row
maximum cosine similarity to other rows is taken as a duplicate score~\cite{salton1988term,ramos2003using}. A row is
flagged when its maximum similarity exceeds a threshold $\tau_{\mathrm{dup}}=0.85$. As Section~\ref{sec:results} shows,
this component is simultaneously the most valuable for detection and the dominant source of false positives, because
shared financial vocabulary inflates pairwise similarity; we analyse and discuss the resulting threshold sensitivity
explicitly.

\subsubsection{Cross-Field Semantic Validation}
Univariate statistics cannot detect inconsistencies that span related fields. Algorithm~\ref{alg:crossfield} implements a
linear-time check for name/identifier alignment that generalises to other field pairs (for example, account-code to
description). The procedure tokenises both fields, ignores tokens of length $\le 2$, and flags a row when no name token
appears as a substring of any identifier token. This check is deterministic and auditable, which is valuable in
regulated contexts even where its marginal effect on aggregate F1 is small (Section~\ref{subsec:ablation}).

\begin{algorithm}[tbp]
\caption{Cross-Field Consistency Check}
\label{alg:crossfield}
\begin{algorithmic}[1]
\Require DataFrame $\mathrm{df}$ with identifier and name columns
\Ensure Set of inconsistent row indices
\State $\mathrm{issues} \gets \emptyset$
\For{each row $r$ in $\mathrm{df}$}
    \State $\mathit{id} \gets \textsc{lower}(r.\mathrm{identifier})$;\ \ $\mathit{nm} \gets \textsc{lower}(r.\mathrm{name})$
    \If{$\mathit{id}$ empty \textbf{or} $\mathit{nm}$ empty} \State \textbf{continue} \EndIf
    \State $I \gets \textsc{tokenize}(\textsc{prefix}(\mathit{id}))$;\ \ $N \gets \textsc{tokenize}(\mathit{nm})$
    \State $\mathit{found} \gets \textsc{False}$
    \For{each $t \in N$ with $|t| > 2$}
        \For{each $e \in I$}
            \If{$t$ is a substring of $e$} \State $\mathit{found} \gets \textsc{True}$ \EndIf
        \EndFor
    \EndFor
    \If{\textbf{not} $\mathit{found}$} \State $\mathrm{issues} \gets \mathrm{issues} \cup \{r\}$ \EndIf
\EndFor
\State \Return $\mathrm{issues}$
\end{algorithmic}
\end{algorithm}

\subsection{Quality-Dimension Scoring}
For reporting, dataset-level quality is summarised across five dimensions. Completeness (Eq.~\ref{eq:completeness}) is the
fraction of populated cells; validity (Eq.~\ref{eq:validity}) the fraction of entries satisfying format and domain
constraints; uniqueness (Eq.~\ref{eq:uniqueness}) the fraction of distinct rows; consistency (Eq.~\ref{eq:consistency}) the
complement of the format-variation rate; and accuracy (Eq.~\ref{eq:accuracy_dim}) a bounded penalty on detected anomalies.

\begin{equation}
\mathrm{Completeness} = \frac{\text{populated cells}}{\text{total cells}} \times 100
\label{eq:completeness}
\end{equation}
\begin{equation}
\mathrm{Validity} = \frac{\text{valid entries}}{\text{total entries}} \times 100
\label{eq:validity}
\end{equation}
\begin{equation}
\mathrm{Uniqueness} = \frac{\text{distinct rows}}{\text{total rows}} \times 100
\label{eq:uniqueness}
\end{equation}
\begin{equation}
\mathrm{Consistency} = \left(1 - \frac{\text{format variations}}{\text{total records}}\right) \times 100
\label{eq:consistency}
\end{equation}
\begin{equation}
\mathrm{Accuracy} = \max\!\left(0,\ 100 - 2\,\times\,\text{anomalies}\right)
\label{eq:accuracy_dim}
\end{equation}

The five dimensions are aggregated with fixed weights (Eq.~\ref{eq:overall}), summing to unity, with completeness and
validity weighted most heavily~\cite{pipino2002understanding,strong1997data}:
\begin{equation}
Q = 0.25\,C + 0.25\,V + 0.20\,U + 0.15\,\mathit{CS} + 0.15\,A
\label{eq:overall}
\end{equation}
where $C,V,U,\mathit{CS},A$ denote completeness, validity, uniqueness, consistency, and accuracy respectively. These
dimensions are descriptive reporting aids; the quantitative detection results in Section~\ref{sec:results} are based on
the row-level confusion matrices, not on $Q$.

\subsection{Composite Confidence and Routing}
\label{subsec:confidence}
Each row receives a composite confidence score that combines the normalised Isolation-Forest anomaly score, the TF--IDF
maximum cosine similarity, and the per-row null rate, as in Eq.~\ref{eq:composite}:
\begin{equation}
s = 0.45\,s_{\mathrm{IF}} + 0.35\,s_{\mathrm{TFIDF}} + 0.20\,s_{\mathrm{null}}
\label{eq:composite}
\end{equation}
with all sub-scores normalised to $[0,1]$. The score is mapped to a routing decision by two thresholds: rows with
$s<\theta_{\mathrm{low}}$ are auto-accepted as likely clean, rows with $s>\theta_{\mathrm{high}}$ are auto-flagged for
rejection, and the remainder are sent to manual review. The default operating point is
$\theta_{\mathrm{low}}=0.40,\ \theta_{\mathrm{high}}=0.70$; a tightened regime ($0.30/0.80$) trades review volume for
automation. This router is what makes the high-recall union strategy practical, and Section~\ref{subsec:confidence_results}
characterises the resulting score distributions and routing rates.

\subsection{Translation Engine and Confidence}
Algorithm~\ref{alg:pipeline} gives the translation pipeline. The source query is analysed for complexity and dialect
features, translated via the LLM, post-processed, and validated; on any error the deterministic rule-based converter is
invoked. A confidence score is then computed (Algorithm~\ref{alg:certainty}) from a base value, a complexity penalty, a
dialect-feature penalty, and syntax/error penalties, and clamped to $[0,1]$. The rule-based path uses a base confidence of
$0.92$; the LLM path can reach higher when validation is clean. The confidence drives the same three-tier routing policy
used for detection, so that low-confidence translations are escalated to human review rather than deployed.

\begin{algorithm}[tbp]
\caption{Confidence-Aware Translation with Fallback}
\label{alg:pipeline}
\begin{algorithmic}[1]
\Require Source query $Q_{\mathrm{src}}$, dialects $D_{\mathrm{src}}, D_{\mathrm{tgt}}$
\Ensure Target query $Q_{\mathrm{tgt}}$, confidence $c$, status
\State $a \gets \textsc{analyse}(Q_{\mathrm{src}})$;\ \ $f \gets \textsc{features}(Q_{\mathrm{src}})$
\State $Q_{\mathrm{tgt}} \gets \textsc{postprocess}(\textsc{llmTranslate}(Q_{\mathrm{src}}, D_{\mathrm{tgt}}))$
\If{translation error} \State $Q_{\mathrm{tgt}} \gets \textsc{ruleFallback}(Q_{\mathrm{src}}, D_{\mathrm{tgt}})$ \EndIf
\State $v \gets \textsc{validate}(Q_{\mathrm{tgt}})$;\ \ $c \gets \textsc{confidence}(a, v)$
\If{$c \ge 0.95$} \State $\mathrm{status} \gets \textsc{auto\_approve}$
\ElsIf{$c \ge 0.70$} \State $\mathrm{status} \gets \textsc{review}$
\Else \State $\mathrm{status} \gets \textsc{manual}$ \EndIf
\State \Return $Q_{\mathrm{tgt}}, c, \mathrm{status}$
\end{algorithmic}
\end{algorithm}

\begin{algorithm}[tbp]
\caption{Translation Confidence Estimation}
\label{alg:certainty}
\begin{algorithmic}[1]
\Require Analysis $a$, validation $v$
\Ensure Confidence $c \in [0,1]$
\State $c \gets 0.92$
\State $c \gets c - \min(0.20,\ 0.01\,a.\mathrm{difficulty})$
\State $c \gets c - 0.05\,|a.\mathrm{features}|$
\If{\textbf{not} $v.\mathrm{syntax\_ok}$} \State $c \gets c - 0.30$ \EndIf
\State $c \gets c - 0.10\,|v.\mathrm{errors}| - 0.02\,|v.\mathrm{warnings}|$
\State \Return $\max(0,\ \min(1,\ c))$
\end{algorithmic}
\end{algorithm}

% =====================================================================
\section{Experimental Setup and Datasets}
\label{sec:setup}

\subsection{Reproducibility Controls}
All experiments are seeded with \texttt{RANDOM\_SEED=42}: global NumPy state, every \texttt{default\_rng} generator, and
every scikit-learn estimator (\texttt{random\_state=42}) are fixed. Datasets are generated once and persisted as static
CSV files; ground-truth masks are reconstructed deterministically from saved JSON. Experiments exercise the analytical
components directly (no HTTP server in the loop), which makes them portable and continuous-integration friendly. Reported
metrics are rounded to four decimals to absorb sub-$10^{-3}$ floating-point differences across library versions; the set
of flagged rows is fully deterministic given the seed. All timing was measured on a single Windows~11 machine (single
process, no parallelism, CPU-only); we therefore treat absolute latencies as conservative upper bounds and discuss
platform sensitivity in Section~\ref{subsec:threats}.

\subsection{Datasets}
\label{subsec:datasets}
Five synthetic financial datasets with known injected-anomaly masks form the labelled corpus, supplemented by a
seven-point scalability series, five noise variants, and a ten-statement SQL suite.
Table~\ref{tab:datasets} lists the labelled datasets, their sizes, and their injected anomaly rates. The anomaly types are
domain-specific: invalid account codes, negative asset balances, and duplicate names for general-ledger accounts;
unbalanced debit/credit pairs and period-format violations for trial balances; same-transaction debit/credit equality and
amount outliers for journal entries; unmapped targets and high-similarity duplicates for the mapping table; and orphaned or
circular nodes for entity structures. We deliberately injected high anomaly rates (\SIrange{10}{38}{\percent}) to stress
the detectors; Section~\ref{subsec:threats} discusses how this differs from production rates and how it affects the
interpretation of precision and FPR.

% TABLE 3 — labelled datasets
\begin{table}[tbp]
\centering
\caption{Labelled Evaluation Datasets (Synthetic, Seed~42)}
\label{tab:datasets}
\small\setlength{\tabcolsep}{4pt}
\begin{tabular}{|l|c|c|c|}
\hline
\textbf{Dataset} & \textbf{Rows} & \textbf{Anomaly Rate} & \textbf{Domain} \\
\hline
GL accounts & 5{,}250 & 22.9\% & General ledger \\
\hline
Journal entries & 8{,}350 & 24.8\% & Debit/credit entries \\
\hline
Trial balance & 3{,}120 & 22.8\% & Period balances \\
\hline
Mapping table & 2{,}000 & 37.5\% & Account-code mapping \\
\hline
Entity structures & 1{,}000 & 10.0\% & Organisational hierarchy \\
\hline
\textbf{Total} & \textbf{19{,}720} & --- & --- \\
\hline
\end{tabular}
\end{table}

\subsection{Metrics and Statistical Protocol}
Classification experiments report precision, recall, F1, FPR, and accuracy from row-level confusion matrices, with
precision and recall defined in Eq.~\ref{eq:pr} and F1 in Eq.~\ref{eq:f1}:
\begin{equation}
\mathrm{Precision} = \frac{\mathrm{TP}}{\mathrm{TP}+\mathrm{FP}}, \qquad
\mathrm{Recall} = \frac{\mathrm{TP}}{\mathrm{TP}+\mathrm{FN}}
\label{eq:pr}
\end{equation}
\begin{equation}
\mathrm{F1} = \frac{2\,\mathrm{Precision}\,\mathrm{Recall}}{\mathrm{Precision}+\mathrm{Recall}}
\label{eq:f1}
\end{equation}
Statistical significance is assessed with a paired two-sample $t$-test ($\alpha=0.05$) over a 10-repeat
$\times$ 5-fold stratified cross-validation (50 paired evaluations), with Cohen's $d$ for effect size and 95\% confidence
intervals computed as $\bar{x}\pm t_{0.025,\,49}\,(s/\sqrt{50})$~\cite{pearson1895notes,han2012data}. Because the
folds are stratified on the anomaly class and the comparison is paired on identical splits, the test controls for
split-level variance~\cite{agrawal1995mining}.

% =====================================================================
\section{Results and Analysis}
\label{sec:results}

\subsection{Detection Quality and Baseline Comparison (RQ1)}
\label{subsec:baseline}
We compared the full hybrid ensemble against six single-method baselines across the four labelled datasets that admit a
shared evaluation (GL accounts, trial balance, journal entries, mapping table). Table~\ref{tab:baseline} reports the mean
metrics, and Figure~\ref{fig:baseline} visualises the same comparison. The hybrid system attains a mean F1 of
$0.321$ at $76.4\%$ recall, substantially exceeding the recall of every principled baseline: rule-based union ($11.3\%$),
Isolation~Forest only ($9.8\%$), and the statistical detectors ($3.3$--$4.4\%$). In financial audit, where a missed
anomaly is far costlier than a redundant review, this recall advantage is the operative quantity.

Two baselines warrant explicit caution. The \emph{string-similarity} baseline shows the highest raw F1 ($0.421$), but it
achieves this through near-universal flagging: its recall is $99.6\%$ \emph{and} its FPR is $99.6\%$, so its apparent F1
is an artefact of the high injected-anomaly rate rather than genuine discrimination, and its $651$\,ms latency makes it
impractical. Conversely, the threshold and IQR detectors post the highest \emph{accuracy} ($\approx71\%$) only because
they predict ``normal'' for almost every row; on imbalanced data, accuracy is misleading and F1/recall are the
appropriate criteria. The hybrid system occupies the balanced operating zone: meaningful F1 and high recall at a
competitive $286$\,ms latency.

% TABLE 4 — baseline comparison (honest numbers)
\begin{table}[tbp]
\centering
\small
\caption{Detector vs.\ Single-Method Baselines (Mean Over 4 Datasets)}
\label{tab:baseline}
\begin{tabular}{|l|c|c|c|c|c|}
\hline
\textbf{System} & \textbf{F1} & \textbf{Prec.} & \textbf{Recall} & \textbf{FPR} & \textbf{ms} \\
\hline
\textbf{Full hybrid} & \textbf{0.321} & 0.264 & \textbf{0.764} & 0.767 & 286 \\
\hline
String similarity$^{\dagger}$ & 0.421 & 0.270 & 0.996 & 0.996 & 651 \\
\hline
TF--IDF only & 0.284 & 0.176 & 0.738 & 0.739 & 142 \\
\hline
Rule-based only & 0.147 & 0.257 & 0.113 & 0.120 & 6 \\
\hline
Isolation Forest only & 0.141 & 0.262 & 0.098 & 0.101 & 122 \\
\hline
Threshold ($z>3$) & 0.072 & 0.214 & 0.044 & 0.043 & 1 \\
\hline
IQR only & 0.052 & 0.127 & 0.033 & 0.030 & 2 \\
\hline
\end{tabular}
\\[2pt]
\raggedright\footnotesize $^{\dagger}$Higher raw F1 is an artefact of near-universal flagging (FPR~$=0.996$); not a
usable detector. See Section~\ref{subsec:baseline}.
\end{table}

\begin{figure}[tbp]
\centering
\includegraphics[width=\columnwidth]{figures/fig2_baseline_comparison.png}
\caption{Baseline comparison across F1, precision, recall, and accuracy. The proposed hybrid system (highlighted)
balances high recall with competitive F1 and latency; the string-similarity bar's high recall coincides with a near-unity
FPR (Table~\ref{tab:baseline}). (New figure, generated from the experiment suite.)}
\label{fig:baseline}
\end{figure}

\subsection{Ablation Study (RQ2)}
\label{subsec:ablation}
We removed each component in turn and measured the change in mean F1 over five datasets. Table~\ref{tab:ablation} and
Figure~\ref{fig:ablation} report the results. TF--IDF semantic matching is decisively the most important component:
removing it drops F1 from $0.312$ to $0.205$ ($\Delta\mathrm{F1}=-0.107$, a $34\%$ relative decline) and collapses recall
from $0.737$ to $0.207$, while also halving latency, confirming it as the principal computational cost. The Isolation
Forest contributes a smaller but real $\Delta\mathrm{F1}=-0.013$, and rule-based validation a marginal
$\Delta\mathrm{F1}=-0.003$ via a recall improvement.

The remaining variants --- removing cross-field validation, confidence routing, or ensemble logic --- show no measurable
change in aggregate F1. We report this honestly rather than suppressing it: binary F1 does not capture the precision/recall
\emph{trade-off} that routing governs, nor the deterministic audit value of cross-field checks on the small subset of rows
they uniquely flag. These components are F1-neutral by construction in this evaluation but operationally meaningful, a
point we return to in Section~\ref{subsec:limitations}.

% TABLE 5 — ablation (honest)
\begin{table}[tbp]
\centering
\caption{Ablation: Mean F1 Over 5 Datasets ($\Delta$ vs.\ Full System)}
\label{tab:ablation}
\begin{tabular}{|l|c|c|c|}
\hline
\textbf{Configuration} & \textbf{F1} & $\boldsymbol{\Delta}$\textbf{F1} & \textbf{Recall} \\
\hline
Full system & 0.312 & --- & 0.737 \\
\hline
w/o Isolation Forest & 0.299 & $-0.013$ & 0.727 \\
\hline
w/o rule-based validation & 0.309 & $-0.003$ & 0.709 \\
\hline
w/o cross-field validation & 0.312 & $0.000$ & 0.737 \\
\hline
w/o confidence routing & 0.312 & $0.000$ & 0.737 \\
\hline
\textbf{w/o TF--IDF (semantic)} & \textbf{0.205} & $\bm{-0.107}$ & \textbf{0.207} \\
\hline
\end{tabular}
\end{table}

\begin{figure}[tbp]
\centering
\includegraphics[width=\columnwidth]{figures/fig1_ablation_study.png}
\caption{Component ablation with F1, precision, and recall per variant and 95\% error bars on F1. Removing TF--IDF
(rightmost groups) causes the steepest degradation; other removals are F1-neutral in this binary evaluation. (New figure,
generated from the experiment suite.)}
\label{fig:ablation}
\end{figure}

\subsection{Statistical Significance}
\label{subsec:stats}
We validated the hybrid system against an Isolation-Forest baseline under 10-repeat $\times$ 5-fold stratified
cross-validation. On GL accounts the full system reached $\mathrm{F1}=0.344\pm0.020$ (95\% CI $[0.338,0.350]$) versus the
baseline's $0.153\pm0.022$ (95\% CI $[0.147,0.159]$), an absolute gain of $+0.191$ (95\% CI of the difference
$[0.183,0.198]$) with $t=51.2$, $p<0.0001$, and Cohen's $d=9.07$. Table~\ref{tab:cv} and Figure~\ref{fig:cv} show the same
pattern across three datasets, with all comparisons significant at $p<0.0001$.

We interpret the very large effect sizes ($d\approx9$--$13$) with care. They do not mean the system is ``nine times
better''; they reflect that the two systems operate in fundamentally different recall regimes ($\approx73\%$ versus
$\approx10\%$) and therefore occupy non-overlapping performance distributions. The fold-level distribution is tight (CV
$=5.8\%$) with no catastrophic fold, supporting the stability of the estimate.

% TABLE 6 — cross validation
\begin{table}[tbp]
\centering
\caption{Cross-Validation: Full System vs.\ Isolation-Forest Baseline}
\label{tab:cv}
\small\setlength{\tabcolsep}{4pt}
\begin{tabular}{|l|c|c|c|c|}
\hline
\textbf{Dataset} & \textbf{Full F1} & \textbf{Base F1} & $\boldsymbol{\Delta}$ & \textbf{Cohen's} $\bm{d}$ \\
\hline
GL accounts & 0.344 & 0.153 & $+0.191$ & 9.1 \\
\hline
Trial balance & 0.371 & 0.140 & $+0.230$ & $\sim$11.8 \\
\hline
Journal entries & 0.395 & 0.138 & $+0.257$ & $\sim$13.2 \\
\hline
\end{tabular}
\\[2pt]
\raggedright\footnotesize All differences significant at $p<0.0001$. $d$ for trial balance and journal entries estimated
from fold distributions.
\end{table}

\begin{figure}[tbp]
\centering
\includegraphics[width=\columnwidth]{figures/fig7_cross_validation_ci.png}
\caption{Cross-validated F1 with 95\% confidence intervals for the full system and the Isolation-Forest baseline; ``$**$''
denotes $p<0.01$. Intervals are tight and non-overlapping across datasets. (New figure, generated from the experiment
suite.)}
\label{fig:cv}
\end{figure}

\subsection{Scalability (RQ3a)}
\label{subsec:scalability}
We profiled the detector across seven dataset sizes from 100 to \num{100000} rows, three timed runs each.
Table~\ref{tab:scalability} and Figure~\ref{fig:scalability} report latency, throughput, and memory. A log--log regression
of latency against size yields an empirical scaling exponent $\alpha=0.69$, i.e.\ sub-linear: per-row cost falls as fixed
overhead amortises. Peak throughput reaches \num{34681} rows/s at \num{100000} rows; memory grows linearly to
$\approx13.8$\,MB. Throughput is non-monotonic, with a dip at \num{50000} rows caused by the $O(n^2)$ cosine-similarity
computation, the system's dominant bottleneck at scale; we mitigate this in production via sparse TF--IDF matrices and
recommend chunked processing beyond \num{50000} rows (Section~\ref{subsec:limitations}).

% TABLE 7 — scalability
\begin{table}[tbp]
\centering
\caption{Scalability (Mean of 3 Runs, Single-Process CPU)}
\label{tab:scalability}
\small\setlength{\tabcolsep}{4pt}
\begin{tabular}{|c|c|c|c|}
\hline
\textbf{Rows} & \textbf{Latency (ms)} & \textbf{Throughput (rows/s)} & \textbf{Mem $\Delta$ (MB)} \\
\hline
100 & 73.6 & 1{,}358 & 0.06 \\
\hline
1{,}000 & 97.5 & 10{,}252 & 0.25 \\
\hline
5{,}000 & 215.2 & 23{,}239 & 0.66 \\
\hline
10{,}000 & 507.9 & 19{,}689 & 1.11 \\
\hline
50{,}000 & 9{,}767.2 & 5{,}119 & 2.52 \\
\hline
100{,}000 & 2{,}883.5 & 34{,}681 & 13.82 \\
\hline
\end{tabular}
\\[2pt]
\raggedright\footnotesize Log--log regression exponent $\alpha=0.69$ (sub-linear). The 50K dip reflects $O(n^2)$ cosine
similarity; see Section~\ref{subsec:scalability}.
\end{table}

\begin{figure}[tbp]
\centering
\includegraphics[width=\columnwidth]{figures/fig3_scalability.png}
\caption{Scalability: (a) latency vs.\ size on log--log axes with fitted slope $\alpha=0.69$; (b) throughput vs.\ size,
showing the 50K cosine-similarity dip; (c) memory growth. (New figure, generated from the experiment suite.)}
\label{fig:scalability}
\end{figure}

\subsection{Robustness Under Noise and Adversarial Inputs (RQ3b)}
\label{subsec:robustness}
We injected five levels of value corruption and nullification and, separately, ran ten pathological edge cases.
Figure~\ref{fig:robustness} summarises both. Recall is stable across all noise levels (\SIrange{93.8}{96.7}{\percent},
under $3$\,pp variance), confirming that the detector does not collapse as data degrades. F1 and precision \emph{rise}
with noise; this is expected for a union strategy on this label model: as more rows become genuinely corrupted, the
high-recall detector's true-positive rate increases and noise itself produces detectable features. The behaviour is
appropriate rather than anomalous, but we flag it explicitly because it would be misleading to present rising F1 as a
quality improvement without this explanation. All ten adversarial cases --- empty frame, all-nulls, single/two rows,
all-identical, mixed dtypes, unicode, oversized strings, extreme numerics, and corrupted headers --- completed without an
unhandled exception, exercising the per-stage fallback paths (zero-return on insufficient sample size, silent rule-skip on
missing columns).

\begin{figure}[tbp]
\centering
\includegraphics[width=\columnwidth]{figures/fig4_robustness.png}
\caption{Robustness: (a) precision/recall/F1 across noise fractions, with recall stable near $0.95$; (b) adversarial
edge-case outcomes (all pass = no unhandled exception). Panel~(b) reports failure-freedom, not detection accuracy. (New
figure, generated from the experiment suite.)}
\label{fig:robustness}
\end{figure}

\subsection{Error and False-Positive Analysis}
\label{subsec:error}
Table~\ref{tab:confusion} gives the per-dataset confusion matrices and Figure~\ref{fig:confusion} their normalised form.
The picture is deliberately unvarnished. On the three text-rich financial datasets the detector achieves near-perfect
recall (FN~$=0$ on trial balance and journal entries; FN~$=3.75\%$ on GL accounts) but at a high FPR ($0.96$--$1.00$),
because TF--IDF cosine similarity accounts for \SIrange{99}{100}{\percent} of false positives: shared financial
terminology (for example, ``accounts payable'') pushes distinct rows above the $0.85$ similarity threshold. Raising the
threshold to $0.90$--$0.92$ or applying domain stop-words would reduce FP substantially at some recall cost. The mapping
table exhibits the inverse pattern --- low FPR ($0.10$) but high FN ($90.7\%$) --- because it lacks text columns for
TF--IDF, leaving the Isolation Forest unable to separate code-only mapping anomalies; this dataset needs code-structure
rules rather than general ML. Crucially, the headline FPR figures reflect both the union strategy and the high synthetic
anomaly rate; under production rates of \SIrange{1}{5}{\percent} and confidence routing, the great majority of these
flags fall into the manual-review tier rather than into rejection (Section~\ref{subsec:confidence_results}).

% TABLE 8 — confusion matrices
\begin{table}[tbp]
\centering
\caption{Per-Dataset Confusion Matrices (Full Hybrid System)}
\label{tab:confusion}
\begin{tabular}{|l|c|c|c|c|c|}
\hline
\textbf{Dataset} & \textbf{TP} & \textbf{FP} & \textbf{FN} & \textbf{Recall} & \textbf{FPR} \\
\hline
GL accounts & 1{,}155 & 3{,}904 & 45 & 0.963 & 0.964 \\
\hline
Trial balance & 710 & 2{,}410 & 0 & 1.000 & 1.000 \\
\hline
Journal entries & 2{,}070 & 6{,}280 & 0 & 1.000 & 1.000 \\
\hline
Mapping table & 70 & 130 & 679 & 0.094 & 0.104 \\
\hline
\end{tabular}
\end{table}

\begin{figure}[tbp]
\centering
\includegraphics[width=\columnwidth]{figures/fig6_confusion_matrices.png}
\caption{Normalised confusion matrices for the four labelled datasets. Text-rich datasets show near-perfect recall with
high FPR (TF--IDF-driven); the mapping table shows the inverse. (New figure, generated from the experiment suite.)}
\label{fig:confusion}
\end{figure}

\subsection{Confidence Score Behaviour and Routing (RQ3c)}
\label{subsec:confidence_results}
Figure~\ref{fig:confidence} shows the composite-score distributions. The three financial datasets share a consistent
right-skewed, leptokurtic profile (mean $\approx0.41$--$0.49$, skewness $\approx1.2$), while the mapping table is shifted
lower (mean $0.10$) because the TF--IDF term is inactive. Under the default thresholds, routing places roughly $13\%$ of
rows in auto-accept, $86\%$ in manual review, and $1\%$ in auto-reject on the financial datasets --- a conservative
profile appropriate to audit. This is precisely the mechanism by which the high-recall, high-FPR detector becomes
operationally usable: rather than auto-rejecting on a single threshold, the system escalates uncertain rows to review and
auto-accepts only confidently clean ones. An F1-maximising threshold sweep independently selects $\approx0.40$ on three
of four datasets, validating the default $\theta_{\mathrm{low}}$.

\begin{figure}[tbp]
\centering
\includegraphics[width=\columnwidth]{figures/fig5_confidence_distributions.png}
\caption{Composite confidence-score distributions per dataset with the F1-optimal threshold marked. Financial datasets are
right-skewed near $0.4$; the mapping table is shifted lower (inactive TF--IDF term). (New figure, generated from the
experiment suite.)}
\label{fig:confidence}
\end{figure}

\subsection{SQL Translation and Confidence (RQ4)}
\label{subsec:translation}
On the ten-statement suite spanning MySQL, PostgreSQL, and SQL~Server sources, the engine produced translations at a mean
confidence of $0.920$ and a mean semantic similarity of $0.968$, with sub-millisecond per-statement latency on the
rule-based path. ANSI-compatible constructs (\texttt{SELECT}, \texttt{JOIN}, \texttt{GROUP BY}, window functions) passed
through unchanged; dialect-specific functions (\texttt{IFNULL}$\rightarrow$\texttt{NVL},
\texttt{AUTO\_INCREMENT}$\rightarrow$\texttt{AUTOINCREMENT}) were correctly transformed. The lowest-fidelity case was a DDL
statement (semantic similarity $0.75$) requiring four structural transforms, consistent with the known difficulty of
\texttt{CREATE TABLE} option syntax. In this run the LLM path was not exercised and all translations used the
deterministic fallback, which is itself an important reliability result: even without the LLM, the system produces
auditable translations at $0.92$ confidence. The per-statement fidelity is shown in Figure~\ref{fig:enterprise}(b).

% =====================================================================
\section{Enterprise Case Study}
\label{sec:casestudy}
To evaluate the integrated pipeline under production-like conditions, we executed a six-stage migration scenario on the
GL-accounts dataset (5{,}250 rows) with the ten-statement SQL suite. Figure~\ref{fig:enterprise}(a) gives the per-stage
latency. End-to-end processing completed in $264.8$\,ms, of which anomaly detection accounted for $94.5\%$ ($250.3$\,ms);
ingestion, account mapping, translation, validation, and reporting together consumed under $14$\,ms. The detector flagged
$96.4\%$ of rows --- expected given the $22.9\%$ injected anomaly rate combined with union logic --- which confidence
routing would tier into $\approx13\%$ auto-accept, $\approx86\%$ review, and $\approx1\%$ reject rather than rejecting
wholesale.

The validation gate correctly \emph{failed} the run, surfacing three issues: an anomaly rate above the $15\%$ threshold,
zero mapping coverage (no target chart-of-accounts was supplied in the experiment), and consequently low mapping
confidence. We report this failure as a positive result: the gate behaved exactly as intended, refusing to certify a
pipeline whose preconditions were unmet, rather than silently passing. SQL translation within the pipeline matched the
standalone results ($0.92$ confidence, $0.968$ mean semantic similarity), with the per-statement fidelity in
Figure~\ref{fig:enterprise}(b). The case study confirms that the dominant cost is detection, that the translation path is
effectively free at this scale, and that the validation gate provides a meaningful safety interlock.

\begin{figure*}[tp]
\centering
\includegraphics[width=0.92\textwidth]{figures/fig8_enterprise_pipeline.png}
\caption{Enterprise case study. (a) Per-stage latency waterfall: anomaly detection dominates at $94.5\%$ of the
$264.8$\,ms total. (b) Per-statement SQL translation semantic similarity grouped by complexity; the DDL statement is the
single lower-fidelity case. (Replaces the previous empty boxplot panel; regenerated from enterprise results.)}
\label{fig:enterprise}
\end{figure*}

% =====================================================================
\section{Security Considerations}
\label{sec:security}
Because the system ingests untrusted tabular data and SQL text, we evaluated twelve adversarial payloads spanning three
severities: SQL injection and XSS and command-injection strings (high); null-byte injection, path traversal, and unicode
adversarial sequences such as right-to-left overrides and zero-width characters (medium); and float overflow, oversized
strings, boundary frames, empty fields, duplicate columns, and extremely wide tables (low). The pass criterion is that the
analytical core processes the input without raising an unhandled exception and without executing any injected content.
Figure~\ref{fig:security} shows that all twelve cases passed. We are careful to frame this as resilience to malformed and
malicious \emph{input handling}, not as a security guarantee for the deployed service: it demonstrates that adversarial
payloads are treated as inert data by the detection path, but production deployments must still enforce authentication,
parameterised database access, output encoding, and transport security at the API boundary. Translated SQL is never
executed by the framework; it is returned for review under the confidence-routing policy, which further limits the blast
radius of an incorrect or malicious translation.

\begin{figure}[tbp]
\centering
\includegraphics[width=\columnwidth]{figures/fig9_security_validation.png}
\caption{Security input-handling validation: twelve adversarial payloads, colour-coded by severity, all processed without
unhandled exceptions. This measures failure-freedom of input handling, not end-to-end deployment security
(Section~\ref{sec:security}). (New figure, generated from the experiment suite.)}
\label{fig:security}
\end{figure}

% =====================================================================
\section{Discussion}
\label{sec:discussion}

\subsection{Interpreting the Operating Point}
The defining characteristic of DataFlow~AI is its deliberate choice of a high-recall operating point governed by a
confidence router, rather than a balanced single-threshold detector. This is appropriate for financial integrity, where
the asymmetric cost of missed anomalies favours recall, but it produces a high false-positive rate that would be
unacceptable if presented to reviewers undifferentiated. The confidence router resolves this tension by escalating only
uncertain rows to review and auto-accepting confidently clean ones, converting raw recall into a manageable workflow. We
have reported the unmanaged FPR and the full confusion matrices precisely so that the trade-off is legible, in contrast to
single-threshold reporting that can mask it.

\subsection{Hybrid Design and Component Roles}
The ablation establishes a clear hierarchy: TF--IDF semantic matching dominates detection but also dominates both false
positives and latency, making its threshold the principal tuning lever; the Isolation Forest contributes secondary lift;
and rules and cross-field checks add deterministic, auditable coverage on specific anomaly types even where their
aggregate-F1 effect is negligible. On the translation side, the rule-based fallback's ability to deliver $0.92$-confidence
translations \emph{without} invoking the LLM is a robustness asset: the system degrades gracefully when the model is
unavailable or rate-limited.

\subsection{Threats to Validity}
\label{subsec:threats}
\emph{Construct validity.} Ground-truth labels are placed synthetically by an independent seeded generator, so
precision/recall measure overlap with a known mask rather than detection of intrinsically anomalous features; this
inflates recall and depresses apparent precision relative to expert-labelled data. The injected anomaly rates
(\SIrange{10}{38}{\percent}) far exceed production rates (\SIrange{1}{5}{\percent}), which raises F1 for recall-oriented
methods and worsens precision; we therefore report FPR and confusion matrices alongside F1 so readers can reason about
other operating points. \emph{Internal validity.} Datasets and baselines were produced by the same team and codebase; we
mitigate alignment risk by placing labels independently of feature generation and by publishing all baseline
implementations. The TF--IDF threshold was set from domain knowledge, not cross-validated, and the error analysis shows
this is the dominant sensitivity. \emph{External validity.} All data are synthetic, single-schema financial records;
generalisation to other ERPs, accounting standards, multi-currency consolidation, and non-financial domains is unverified.
Scalability was measured to \num{100000} rows on a single CPU; million-row and multi-core behaviour is projected, not
measured. SQL coverage is limited to three source dialects into Snowflake/PostgreSQL; Oracle PL/SQL, Teradata, and
BigQuery constructs are untested. \emph{Statistical validity.} The very large effect sizes reflect non-overlapping recall
regimes between systems and should be read as ``qualitatively different operating points'', not multiplicative
superiority. Table~\ref{tab:threats} summarises these threats and their status.

% TABLE 9 — threats to validity
\begin{table}[tbp]
\centering
\caption{Threats to Validity and Mitigations}
\label{tab:threats}
\begin{tabular}{|l|c|c|}
\hline
\textbf{Threat} & \textbf{Type} & \textbf{Severity} \\
\hline
Synthetic ground truth & Construct & High \\
\hline
High injected anomaly rate & Construct & Medium \\
\hline
In-house baselines & Internal & Medium \\
\hline
Un-tuned TF--IDF threshold & Internal & Medium \\
\hline
Single financial domain & External & High \\
\hline
100K-row scalability ceiling & External & Medium \\
\hline
Limited SQL dialect coverage & External & Medium \\
\hline
Single-platform timing & External & Low \\
\hline
Effect-size magnitude & Statistical & Low \\
\hline
\end{tabular}
\end{table}

\subsection{Reproducibility}
\label{subsec:repro}
Every experiment is governed by a single seed (42) across NumPy, generators, and estimators; datasets are generated once
and read-only thereafter; ground-truth masks are reconstructed deterministically; and outputs are written as ordered JSON.
Expected checksums (for example, full-system ablation $\mathrm{F1}=0.3119$, no-TF--IDF $\mathrm{F1}=0.2054$, baseline
hybrid $\mathrm{F1}=0.3212$, scalability $\alpha=0.69$, CV Cohen's $d=9.07$) allow third parties to verify a re-run to
within $10^{-3}$. Pinned dependencies and an ordered master runner are provided. The only non-determinism is sub-$10^{-3}$
floating-point variation across library versions, absorbed by four-decimal rounding.

\subsection{Limitations and Future Work}
\label{subsec:limitations}
The principal limitations follow directly from the threats above. The fixed TF--IDF threshold is the dominant precision
lever and should be replaced by cross-validated, domain-specific selection. Cross-field validation currently covers
name/identifier pairs and requires extension for additional relations. Datasets beyond \num{50000} rows incur the
$O(n^2)$ cosine-similarity cost and should use chunked processing; distributed execution (e.g.\ Spark or Dask) is required
for million-row scale. Translation coverage should extend to Oracle, Teradata, and BigQuery, with formal semantic
verification rather than token-level similarity. Most importantly, the synthetic evaluation must be complemented by
expert-annotated financial benchmarks and a controlled real-world deployment in which audit outcomes supply true labels.
Near-term work targets adaptive thresholds, self-hosted LLMs, and real-time monitoring; medium-term work targets
distributed processing and active learning from reviewer feedback; longer-term work targets multi-domain validation beyond
finance.

% =====================================================================
\section{Conclusion}
\label{sec:conclusion}
We presented DataFlow~AI, a hybrid, confidence-aware framework that integrates data-integrity validation and cross-dialect
SQL translation in one deployment-oriented pipeline. Under a fully seeded, reproducible evaluation, the hybrid ensemble
achieved a mean F1 of $0.32$ at $76\%$ recall --- a statistically significant $+0.19$ absolute F1 gain over an
Isolation-Forest baseline ($t=51.2$, $p<0.0001$, $d=9.1$) --- with TF--IDF semantic matching identified as the dominant
component ($\Delta\mathrm{F1}=-0.107$ on removal). The system scaled sub-linearly ($\alpha=0.69$), maintained stable recall
under noise, processed all adversarial and security inputs without failure, and translated SQL at $0.92$ confidence and
$0.97$ semantic similarity even on the deterministic fallback path. Rather than reporting a single favourable threshold,
we characterised the high-recall, high-FPR operating point in full and showed how confidence routing converts it into a
practical, audit-appropriate review workflow. The contribution is therefore not a new accuracy record but a transparent,
reproducible demonstration that hybrid, confidence-aware integration of validation and translation is both feasible and
operationally sensible --- subject to the clearly stated limits of synthetic evaluation, which expert-labelled and
real-world studies should next address.

% =====================================================================
\begin{thebibliography}{99}

\bibitem{redman1998impact}
T. Redman, ``The impact of poor data quality on the typical enterprise,'' \textit{Commun. ACM}, vol. 41, no. 2, pp. 79--82, 1998.

\bibitem{hultgren2018data}
G. Hultgren, \textit{Data Quality: Why It Matters and How to Improve It}. O'Reilly Media, 2018.

\bibitem{stonebraker2007end}
M. Stonebraker et al., ``The end of an architectural era: (It's time for a complete rewrite),'' in \textit{Proc. VLDB}, 2007, pp. 1150--1160.

\bibitem{batini2009data}
C. Batini, C. Cappiello, C. Francalanci, and A. Maurino, ``Methodologies for data quality assessment and improvement,'' \textit{ACM Comput. Surv.}, vol. 41, no. 3, pp. 16:1--16:52, 2009.

\bibitem{ilyas2015functional}
I. F. Ilyas and X. Chu, ``Trends in cleaning relational data: Consistency and other flavors,'' \textit{SIGMOD Rec.}, vol. 44, no. 4, pp. 35--45, 2015.

\bibitem{motro1989integrity}
A. Motro, ``Integrity = Validity + Completeness,'' \textit{ACM Trans. Database Syst.}, vol. 14, no. 4, pp. 480--502, 1989.

\bibitem{sculley2015hidden}
D. Sculley et al., ``Hidden technical debt in machine learning systems,'' in \textit{Proc. NeurIPS}, 2015.

\bibitem{great_expectations}
Great Expectations, ``Great Expectations: Always know what to expect from your data,'' 2023. [Online]. Available: https://greatexpectations.io

\bibitem{pedregosa2011scikit}
F. Pedregosa et al., ``Scikit-learn: Machine learning in Python,'' \textit{J. Mach. Learn. Res.}, vol. 12, pp. 2825--2830, 2011.

\bibitem{aggarwal2017outlier}
C. C. Aggarwal, \textit{Outlier Analysis}, 2nd ed. Springer, 2017.

\bibitem{rahm2001survey}
E. Rahm and P. A. Bernstein, ``A survey of approaches to automatic schema matching,'' \textit{VLDB J.}, vol. 10, pp. 334--350, 2001.

\bibitem{manning2008introduction}
C. D. Manning, P. Raghavan, and M. Sch\"{u}tze, \textit{Introduction to Information Retrieval}. Cambridge Univ. Press, 2008.

\bibitem{ribeiro2016why}
M. Ribeiro, S. Singh, and C. Guestrin, ``Why should I trust you? Explaining the predictions of any classifier,'' in \textit{Proc. KDD}, 2016, pp. 1135--1144.

\bibitem{neumann1996risks}
P. G. Neumann, ``Risks of automation,'' \textit{Commun. ACM}, vol. 39, no. 3, pp. 150--152, 1996.

\bibitem{lundberg2017unified}
S. M. Lundberg and S.-I. Lee, ``A unified approach to interpreting model predictions,'' in \textit{Proc. NeurIPS}, 2017, pp. 4765--4774.

\bibitem{liu2008isolation}
F. T. Liu, K. M. Ting, and Z.-H. Zhou, ``Isolation Forest,'' in \textit{Proc. IEEE ICDM}, 2008, pp. 413--422.

\bibitem{salton1988term}
G. Salton and C. Buckley, ``Term-weighting approaches in automatic text retrieval,'' \textit{Inf. Process. Manage.}, vol. 24, no. 5, pp. 513--523, 1988.

\bibitem{ramos2003using}
J. Ramos, ``Using TF-IDF to determine word relevance in document queries,'' in \textit{Proc. First Instructional Conf. Machine Learning}, 2003.

\bibitem{pipino2002understanding}
L. L. Pipino, Y. W. Lee, and R. Y. Wang, ``Data quality assessment,'' \textit{Commun. ACM}, vol. 45, no. 4, pp. 211--218, 2002.

\bibitem{strong1997data}
D. M. Strong, Y. W. Lee, and R. Y. Wang, ``Data quality in context,'' \textit{Commun. ACM}, vol. 40, no. 5, pp. 103--110, 1997.

\bibitem{pearson1895notes}
K. Pearson, ``Notes on regression and inheritance in the case of two parents,'' \textit{Proc. Royal Soc. London}, vol. 58, pp. 240--242, 1895.

\bibitem{han2012data}
J. Han, M. Kamber, and J. Pei, \textit{Data Mining: Concepts and Techniques}, 3rd ed. Morgan Kaufmann, 2012.

\bibitem{agrawal1995mining}
R. Agrawal and R. Srikant, ``Mining sequential patterns,'' in \textit{Proc. ICDE}, 1995, pp. 3--14.

\end{thebibliography}

% =====================================================================
\begin{IEEEbiography}[{
\includegraphics[width=1in,height=1.25in,clip,keepaspectratio]{figures/author_upasana.png}}]{Upasana Bhaumik}
is an undergraduate student in Computer Science and Engineering with a specialization in Bioinformatics at Vellore
Institute of Technology, India. She builds AI-driven and data-centric applications, combining machine learning, data
engineering, and full-stack development. Her interests include large-language-model-powered systems, data-quality
automation, and applications of AI in bioinformatics and healthcare.
\end{IEEEbiography}

\begin{IEEEbiography}[{\includegraphics[width=1in,height=1.25in,clip,keepaspectratio]{figures/author_berlin.png}}]{Berlin M A}
received the master's degree in computer science engineering from National Engineering College, Tamil Nadu, and the
Ph.D. degree in computer science engineering from Anna University, Chennai. She is currently an Associate Professor with
the School of Computer Science and Engineering, Vellore Institute of Technology, Vellore. Her research interests include
ad hoc networks, the Internet of Things, and machine learning. She has over 20 years of teaching and research experience,
with numerous publications and patents in these areas.
\end{IEEEbiography}

\EOD

\end{document}